% concatenated from main.tex
\documentclass[11pt]{article}

\usepackage[letterpaper,margin=1in]{geometry}
\usepackage[T1]{fontenc}
\usepackage[utf8]{inputenc}
\usepackage{amsmath,amssymb,amsthm,mathtools}
\usepackage{newpxtext}
\usepackage{newpxmath}
\usepackage{microtype}
\usepackage[dvipsnames]{xcolor}
\usepackage{aliascnt}
\usepackage[hypertexnames=false]{hyperref}
\usepackage[nameinlink,noabbrev]{cleveref}

\hypersetup{
  colorlinks=true,
  linkcolor=NavyBlue,
  citecolor=ForestGreen,
  urlcolor=NavyBlue,
  pdftitle={The Balancing Theorem for Spanning Trees of Rectangular Grid Graphs},
  pdfauthor={}
}

\numberwithin{equation}{section}

\newtheorem{theorem}{Theorem}[section]

\newaliascnt{corollary}{theorem}
\newtheorem{corollary}[corollary]{Corollary}
\aliascntresetthe{corollary}

\newaliascnt{conjecture}{theorem}
\newtheorem{conjecture}[conjecture]{Conjecture}
\aliascntresetthe{conjecture}

\newaliascnt{lemma}{theorem}
\newtheorem{lemma}[lemma]{Lemma}
\aliascntresetthe{lemma}

\newaliascnt{proposition}{theorem}
\newtheorem{proposition}[proposition]{Proposition}
\aliascntresetthe{proposition}

\crefname{theorem}{Theorem}{Theorems}
\crefname{corollary}{Corollary}{Corollaries}
\crefname{conjecture}{Conjecture}{Conjectures}
\crefname{lemma}{Lemma}{Lemmas}
\crefname{proposition}{Proposition}{Propositions}
\Crefname{theorem}{Theorem}{Theorems}
\Crefname{corollary}{Corollary}{Corollaries}
\Crefname{conjecture}{Conjecture}{Conjectures}
\Crefname{lemma}{Lemma}{Lemmas}
\Crefname{proposition}{Proposition}{Propositions}

\DeclareMathOperator{\arcosh}{arcosh}
\DeclareMathOperator{\arsinh}{arsinh}
\newcommand{\Z}{\mathbb{Z}}
\newcommand{\dd}{\,d}

\title{\bfseries The Balancing Theorem for Spanning Trees of Rectangular Grid Graphs}
\author{Jiechen Zhang\footnote{EPFL\@. Email: \href{mailto:jiechen.zhang@epfl.ch}{jiechen.zhang@epfl.ch}.}}
\date{}

\begin{document}
\maketitle

\begin{abstract}
We compare spanning-tree counts of rectangular grid graphs at fixed area. The main result is a balancing theorem: if $AB=ab$ and $A\le a\le b\le B$, then the more balanced rectangle has larger spanning-tree count, and in fact
\[
  \log\tau(a,b)-\log\tau(A,B)
  \ge
  \operatorname{arsinh}(1)\bigl((A+B)-(a+b)\bigr).
\]
The universal constant $\operatorname{arsinh}(1)=\log(1+\sqrt{2})\approx0.88137$ is optimal. Thus the log-ratio between equal-area rectangles grows at least linearly in the reduction of the sum of side lengths. The proof starts from the Laplacian product formula, passes to hyperbolic coordinates, and gives an exact finite-size identity for the log-gain of a balancing move. A monotone trapezoidal estimate gives the optimal main term, while a positive monotone residual gives strictness. The same decomposition identifies the closest-to-square divisor rectangle as the unique fixed-area maximizer and recovers the bounded-aspect rectangular asymptotic through the logarithmic corner term, including the square-grid case.
\end{abstract}

\clearpage

\section{Introduction}

Let $P_r$ be the path graph on $r$ vertices, and let $R_{r,s}=P_r\square P_s$ be the $r\times s$ rectangular grid graph formed by their Cartesian product. Let $\tau(G)$ denote the number of spanning trees of a finite connected graph $G$, and abbreviate $\tau(r,s)=\tau(R_{r,s})$.

Spanning-tree counts are a classical graph invariant, and rectangular grids are among the basic finite pieces of the square lattice for which product formulas make exact finite-size questions accessible. The rectangular case isolates a natural finite-size comparison problem: among equal-area rectangles, does the more balanced shape have the larger spanning-tree count?

This question also has a motivation in political redistricting: in tree-weighted graph partition models, weights proportional to spanning-tree counts serve as a compactness-sensitive local statistic.

Procaccia and Tucker-Foltz record a broader square-maximality conjecture: among grid subgraphs with a fixed number of vertices, square subgrids should maximize the spanning-tree count~\cite{ProcacciaTuckerFoltz2022}. In the connected induced-subgraph formulation, let $\mathcal L$ be the infinite square-lattice graph with vertex set $\Z^2$.

\begin{conjecture}\label{conj:induced-grid}
Let $n\ge 1$ be an integer. Let $S\subset\Z^2$ be a finite set of vertices such that $|S|=n^2$, and suppose that the induced subgraph $\mathcal L[S]$ is connected. Then
\[
  \tau(\mathcal L[S])\le \tau(n,n).
\]
\end{conjecture}

We affirmatively answer the rectangular form of this question at the finite level. More precisely, for all equal-area rectangular grid graphs, balancing the side lengths increases the spanning-tree count, and we quantify the gain with an explicit optimal lower bound.

\begin{theorem}[Balancing theorem]\label{thm:balancing}
Let $A,B,a,b$ be positive integers satisfying
\[
  AB=ab,\qquad A\le a\le b\le B.
\]
Then
\[
  \log\tau(a,b)-\log\tau(A,B)
  \ge
  \arsinh(1)\bigl((A+B)-(a+b)\bigr).
\]
If the rectangles are distinct, the inequality is strict. The universal constant $\arsinh(1)=\log(1+\sqrt2)\approx0.88137$ is optimal.
\end{theorem}

The quantity $(A+B)-(a+b)$ is the reduction in the sum of side lengths when the rectangle is balanced. For rectangles embedded in the square lattice, it is also half the reduction in the outer grid-boundary incidence count. The statement is pairwise and exact within the rectangular family, in contrast with coarser perimeter comparisons that do not resolve which equal-area rectangle has larger spanning-tree count.

\subsection{Our Contributions}

We give a finite-size comparison theory for spanning-tree counts of rectangular grid graphs at fixed area. The main contributions are the following.

\begin{enumerate}
\setlength{\itemsep}{1pt}
\setlength{\parsep}{0pt}
\setlength{\topsep}{2pt}
        \item[(1)] \emph{Balancing theorem} (\textbf{\Cref{thm:balancing}}). Among equal-area rectangular grids, moving toward a more balanced shape strictly increases the number of spanning trees, with the optimal universal constant $\alpha=\arsinh(1)$. The proof reduces Kirchhoff's product formula to a one-dimensional hyperbolic identity (\textbf{\Cref{prop:product,lem:hyperbolic,thm:exact-gain}}), then combines a sharp trapezoidal estimate with an optimality sequence (\textbf{\Cref{lem:trap,prop:optimality}}).
\item[(2)] \emph{Divisor-rectangle maximality} (\textbf{\Cref{cor:square-max}}). Among fixed-area divisor rectangles, the closest-to-square divisor pair is the unique maximizer; in particular, the $n\times n$ square is the unique maximizing rectangle at area $n^2$.
\item[(3)] \emph{Residual asymptotics} (\textbf{\Cref{prop:bounded-aspect-residual}}). The residual term is not only positive for the comparison theorem; it also contains the bounded-aspect rectangular asymptotic, including the classical square-grid area, boundary, and logarithmic corner terms.
\end{enumerate}

The comparison is not a factor-by-factor consequence of the product formula. After the Kirchhoff normalization and the one-dimensional modes cancel, the remaining product for $R_{r,s}$ has $(r-1)(s-1)$ interaction factors. Under a fixed-area balancing move, both the number of factors and the eigenvalue grids change. Some spectral quantities increase and others decrease, and there is no natural pointwise matching of factors that proves the result. The proof instead uses the Chebyshev-hyperbolic form of the one-dimensional characteristic polynomial to express the entire log-ratio as a structured sum. Only after splitting
\[
  \log\frac{\sinh(tz)}{\sinh z}
  =
  (t-1)z+\log\frac{1-e^{-2tz}}{1-e^{-2z}}
\]
do the two mechanisms become visible: a main term governed by a trapezoidal identity for the special function $c(x)=\arcosh(2-\cos \pi x)$, and a positive monotone residual.
As a short illustration, the same residual decomposition recovers the correct area, boundary, and logarithmic corner terms in the bounded-aspect rectangular asymptotic, including the square case. The hyperbolic form also gives one-dimensional evaluation formulas for related product grids, although the sharp comparison theorem is specific to rectangles.
For comparison with a more elementary concavity viewpoint, \Cref{app:concavity-constant} records a weaker balancing constant $\alpha_0=c(1)-\int_0^1 c(x)\,dx$ and its proof; the sharp constant in the main theorem remains $\alpha=\arsinh(1)$.

All statements below are restricted to rectangles. We do not analyze arbitrary nonrectangular grid shapes. Within the rectangular class, however, the results are exact: the closest-to-square divisor pair is the unique maximizing shape, and every pairwise rectangular comparison is controlled by the same one-dimensional formula.

\subsection{Related Work}

The classical starting point is Kirchhoff's Matrix--Tree Theorem~\cite{Kirchhoff1847}, which expresses the spanning-tree count as a Laplacian cofactor, together with product formulas and asymptotic studies for rectangular grids and lattice graphs~\cite{Temperley1974,Wu1977,ShrockWu2000,GolinLeungWangYong2005}. Periodic boundary conditions and toroidal grids have also been studied from a spectral-asymptotic viewpoint: Chinta, Jorgenson, and Karlsson related spanning-tree complexity of discrete tori to heat kernels and heights of real tori~\cite{ChintaJorgensonKarlsson2010,ChintaJorgensonKarlsson2012}, and Louis gave factor-reduction formulas for circulant graphs and discrete tori~\cite{Louis2015}.

A separate line of motivation comes from redistricting and tree-weighted graph partitions. ReCom-type chains and reversible variants use these weights~\cite{DeFordDuchinSolomon2021,CannonDuchinRandallRule2026}. Procaccia and Tucker-Foltz relate them to cut-edge perimeter and record the broader square-maximality conjecture for induced grid subgraphs~\cite{ProcacciaTuckerFoltz2022}; Cannon, Pegden, and Tucker-Foltz study balanced tree-weighted grid partitions~\cite{CannonPegdenTuckerFoltz2024}, and Cannon, Pankow, Pegden, and Tucker-Foltz later study direct sampling of tree-weighted partitions~\cite{CannonPankowPegdenTuckerFoltz2025}. Tapp's boundary-sensitive bounds for grid subgraphs are close in spirit~\cite{Tapp2024}. These results support the compactness and square-maximality heuristic, but they do not imply the exact finite equal-area rectangular comparison proved here, even within the rectangular subclass.

\paragraph{Organization.}
\Cref{sec:preliminaries} records the rectangular-grid product formula. \Cref{sec:hyperbolic-setup} introduces the hyperbolic reduction. \Cref{sec:balancing-gain} derives the balancing-gain identity. \Cref{sec:trapezoidal-estimate} proves the trapezoidal estimate that supplies the sharp constant. \Cref{sec:proof-optimality} proves the balancing theorem and optimality, derives the bounded-aspect asymptotic, and gives the divisor-rectangle consequence. \Cref{sec:discussion} discusses the scope of the method and the induced-grid conjecture.

\section{Preliminaries}\label{sec:preliminaries}

For $r\ge 1$ and $0\le j\le r-1$, set
\[
  \lambda_j(r)=2-2\cos\frac{\pi j}{r}.
\]
Thus $\lambda_0(r)=0$, and for $j\ge1$, $\lambda_j(r)=4\sin^2(\pi j/(2r))$.

\begin{proposition}\label{prop:product}
For all positive integers $r,s$,
\[
  \tau(r,s)
  =
  \prod_{j=1}^{r-1}\prod_{k=1}^{s-1}
  \bigl(\lambda_j(r)+\lambda_k(s)\bigr),
\]
with the convention that an empty product is $1$.
\end{proposition}

\begin{proof}
The Laplacian eigenvalues of $P_r$ are $\lambda_j(r)$, $0\le j\le r-1$. Hence the Laplacian eigenvalues of $R_{r,s}=P_r\square P_s$ are $\lambda_j(r)+\lambda_k(s)$, with $0\le j\le r-1$ and $0\le k\le s-1$.
The only zero eigenvalue is the one with $(j,k)=(0,0)$. Kirchhoff's Matrix--Tree Theorem~\cite{Kirchhoff1847} gives
\[
  \tau(r,s)
  =
  \frac{1}{rs}
  \prod_{\substack{0\le j\le r-1\\0\le k\le s-1\\(j,k)\ne(0,0)}}
  \bigl(\lambda_j(r)+\lambda_k(s)\bigr).
\]
The one-dimensional factors cancel the prefactor because $\prod_{j=1}^{r-1}\lambda_j(r)=r$ for every $r\ge1$, with the empty-product convention when $r=1$. For $r\ge2$ this is the standard sine product $\prod_{j=1}^{r-1}2\sin(\pi j/(2r))=\sqrt r$.
After canceling the factors with $j=0$ or $k=0$, only the positive-mode interaction product remains.
\end{proof}

The empty-product convention gives $\tau(1,s)=\tau(r,1)=1$, as it should for one-dimensional path graphs.

\section{Hyperbolic Setup}\label{sec:hyperbolic-setup}

For $r\ge1$, define
\[
  q_r(x)=\prod_{j=1}^{r-1}\bigl(x+\lambda_j(r)\bigr).
\]
Then \Cref{prop:product} can be written as
\begin{equation}\label{eq:tau-q}
  \tau(r,s)
  =
  \prod_{j=1}^{r-1}q_s(\lambda_j(r))
  =
  \prod_{k=1}^{s-1}q_r(\lambda_k(s)).
\end{equation}

\begin{lemma}\label{lem:hyperbolic}
If $x=2\cosh z-2$ with $z>0$, then
\[
  q_r(x)=\frac{\sinh(rz)}{\sinh z}.
\]
\end{lemma}

\begin{proof}
Let $p_1(x)=1$, $p_2(x)=x+2$, and
\[
  p_{r+1}(x)=(x+2)p_r(x)-p_{r-1}(x).
\]
By induction, $p_r$ is monic of degree $r-1$. If $x=2\cosh z-2$, then
\[
  \sinh((r+1)z)=2\cosh z\,\sinh(rz)-\sinh((r-1)z)
\]
shows that $p_r(x)=\sinh(rz)/\sinh z$.

For $r\ge2$, substituting $x=2\cos\phi-2$ gives
\[
  p_r(2\cos\phi-2)=\frac{\sin(r\phi)}{\sin\phi}
\]
whenever $\sin\phi\ne0$. Taking $\phi=\pi j/r$, $1\le j\le r-1$, shows that $p_r(-\lambda_j(r))=0$. Thus $p_r$ and $q_r$ are monic polynomials of degree $r-1$ with the same roots.
\end{proof}

Define
\[
  c(x)=\arcosh(2-\cos\pi x)
  =
  2\arsinh\left(\sin\frac{\pi x}{2}\right),
  \qquad 0\le x\le1.
\]
For $1\le j\le r-1$,
\begin{equation}\label{eq:lambda-c}
  \lambda_j(r)=2\cosh c(j/r)-2.
\end{equation}

\begin{lemma}\label{lem:c}
The function $c$ satisfies $c(0)=0$, is increasing and concave on $[0,1]$, and
\[
  c(1)=2\arsinh(1)=2\log(1+\sqrt{2}).
\]
\end{lemma}

\begin{proof}
The value $c(0)=0$ and the formula for $c(1)$ are immediate from
$c(x)=2\arsinh(\sin(\pi x/2))$.
Differentiating gives
\[
  c'(x)=\frac{\pi\cos(\pi x/2)}{(1+\sin^2(\pi x/2))^{1/2}}\ge0
\]
and
\[
  c''(x)=
  -\frac{\pi^2\sin(\pi x/2)}
  {(1+\sin^2(\pi x/2))^{3/2}}
  \le0.
\]
\end{proof}

Put
\[
  C_r=\sum_{j=1}^{r-1}c(j/r).
\]
We also write
\[
  \alpha=\frac{c(1)}{2}=\arsinh(1)=\log(1+\sqrt{2}).
\]
Finally set
\[
  I=\int_0^1c(x)\dd x.
\]

\section{Balancing Gain}\label{sec:balancing-gain}

We first record an exact identity for the log-gain under a balancing move.

\begin{lemma}\label{lem:G}
Let $t>1$, and define
\[
  G_t(z)=\log\frac{1-e^{-2tz}}{1-e^{-2z}},
  \qquad z>0.
\]
Then $G_t(z)>0$ for all $z>0$, and $G_t$ is strictly decreasing on $(0,\infty)$.
\end{lemma}

\begin{proof}
Since $t>1$, we have $1-e^{-2tz}>1-e^{-2z}$ for every $z>0$, so $G_t(z)>0$. Differentiating gives
\[
  G_t'(z)=\frac{2t}{e^{2tz}-1}-\frac{2}{e^{2z}-1}.
\]
For $y>0$, the function $(e^y-1)/y$ is strictly increasing, since
\[
  \frac{e^y-1}{y}=\int_0^1e^{sy}\dd s.
\]
Thus
\[
  \frac{e^{2tz}-1}{2tz}>\frac{e^{2z}-1}{2z},
\]
which is equivalent to $G_t'(z)<0$.
\end{proof}

\begin{theorem}\label{thm:exact-gain}
Let $A,B,a,b$ be positive integers satisfying
\[
  AB=ab,\qquad A\le a\le b\le B.
\]
Set
\[
  t=\frac{a}{A}=\frac{B}{b}.
\]
If $t=1$, then the rectangles are identical and
\[
  \log\tau(a,b)-\log\tau(A,B)=0.
\]
If $t>1$, then
\begin{equation}\label{eq:exact-gain}
  \log\frac{\tau(a,b)}{\tau(A,B)}
  =
  (t-1)(AC_b-bC_A)+\Gamma,
\end{equation}
where
\[
  \Gamma
  =
  \sum_{k=1}^{b-1}G_t\bigl(Ac(k/b)\bigr)
  -
  \sum_{j=1}^{A-1}G_t\bigl(bc(j/A)\bigr).
\]
Moreover, $\Gamma>0$ when $t>1$.
\end{theorem}

\begin{proof}
The case $t=1$ is immediate. Assume $t>1$. We compare both rectangles to the intermediate rectangle $A\times b$. By \eqref{eq:tau-q}, \eqref{eq:lambda-c}, and \Cref{lem:hyperbolic},
\[
  \log\frac{\tau(a,b)}{\tau(A,b)}
  =
  \sum_{k=1}^{b-1}
  \log\frac{\sinh(a\,c(k/b))}{\sinh(A\,c(k/b))}
\]
and
\[
  \log\frac{\tau(A,B)}{\tau(A,b)}
  =
  \sum_{j=1}^{A-1}
  \log\frac{\sinh(B\,c(j/A))}{\sinh(b\,c(j/A))}.
\]
Subtracting and using $a=tA$ and $B=tb$ gives
\[
  \log\frac{\tau(a,b)}{\tau(A,B)}
  =
  \sum_{k=1}^{b-1}
  \log\frac{\sinh(tA\,c(k/b))}{\sinh(A\,c(k/b))}
  -
  \sum_{j=1}^{A-1}
  \log\frac{\sinh(tb\,c(j/A))}{\sinh(b\,c(j/A))}.
\]
For $z>0$,
\[
  \log\frac{\sinh(tz)}{\sinh z}
  =
  (t-1)z+
  \log\frac{1-e^{-2tz}}{1-e^{-2z}}
  =
  (t-1)z+G_t(z).
\]
This proves \eqref{eq:exact-gain}.

It remains to show that $\Gamma>0$. Since $t>1$, we have $a=tA>A$. Combined with $a\le b$, this implies $b>A$. If $A=1$, there are no terms in the second sum to match. Otherwise, for each $j=1,\ldots,A-1$, set $k_j=\lfloor jb/A\rfloor$.
For these indices, $b/A\le jb/A<b$, so $1\le k_j\le b-1$. If $j'<j$, then $(j-j')b/A\ge b/A>1$, so $k_{j'}<k_j$. Thus the $k_j$ are distinct elements of $\{1,\ldots,b-1\}$. Also $k_j/b\le j/A$. Since $c$ is increasing,
\[
  A c(k_j/b)\le A c(j/A)\le b c(j/A).
\]
By \Cref{lem:G}, $G_t$ is decreasing, and therefore
\[
  G_t\bigl(Ac(k_j/b)\bigr)\ge G_t\bigl(bc(j/A)\bigr).
\]
After matching these $A-1$ terms, the first sum has $b-A$ unmatched terms, each strictly positive by \Cref{lem:G}. Hence $\Gamma>0$.
\end{proof}

\section{A Trapezoidal Estimate}\label{sec:trapezoidal-estimate}

The optimal constant comes from a special trapezoidal monotonicity property of $c$. Let
\[
  \gamma(\theta)=\arcosh(2-\cos\theta),
  \qquad 0\le\theta\le\pi,
\]
so that $c(x)=\gamma(\pi x)$ and
\[
  I=\frac1\pi\int_0^\pi\gamma(\theta)\dd\theta.
\]
For $r\ge1$, define the half-weighted trapezoidal sum
\[
  T_r=
  \frac1r\left(
  \frac{c(0)+c(1)}2+\sum_{j=1}^{r-1}c(j/r)
  \right)
  =
  \frac{C_r}{r}+\frac{\alpha}{r}.
\]
This is the composite trapezoidal approximation to $I$ after the change of variables $\theta=\pi x$. The endpoint contribution is exactly $c(1)/(2r)=\alpha/r$, which is the source of the constant in the balancing theorem. The lemma below gives an exact formula for $T_r-I$; that error is negative and increases monotonically to $0$ as $r$ grows.

\begin{lemma}\label{lem:trap}
For every $r\ge1$,
\begin{equation}\label{eq:trap-identity}
  T_r
  =
  I+
  \frac{1}{\pi r}
  \int_0^\pi
  \log\left(1-e^{-2r\gamma(\phi)}\right)\dd\phi.
\end{equation}
Consequently, if $1\le r\le s$, then
\begin{equation}\label{eq:C-trapezoidal}
  \frac{C_s}{s}-\frac{C_r}{r}
  \ge
  \frac{s-r}{rs}\alpha.
\end{equation}
\end{lemma}

\begin{proof}
We first prove \eqref{eq:trap-identity}. We use the standard identity, valid for $a\ge1$,
\[
  \frac1\pi\int_0^\pi\log(a-\cos\phi)\dd\phi
  =
  \log\frac{a+\sqrt{a^2-1}}{2}.
\]
With $a=2-\cos\theta$, this gives
\begin{equation}\label{eq:gamma-integral}
  \gamma(\theta)
  =
  \log2+
  \frac1\pi\int_0^\pi
  \log(2-\cos\theta-\cos\phi)\dd\phi.
\end{equation}

Fix $0<\phi\le\pi$ and put $q=e^{-\gamma(\phi)}$. Since
\[
  2-\cos\phi=\frac{q+q^{-1}}2,
\]
we have
\[
  2-\cos\theta-\cos\phi
  =
  \frac{1-2q\cos\theta+q^2}{2q}.
\]
Thus, with $L(\theta,\phi)=\log(2-\cos\theta-\cos\phi)$,
\[
  L(\theta,\phi)
  =
  \gamma(\phi)-\log2+
  \log(1-2q\cos\theta+q^2).
\]
For $0<q<1$,
\[
  \log(1-2q\cos\theta+q^2)
  =
  -2\sum_{m=1}^\infty\frac{q^m}{m}\cos(m\theta).
\]
The half-weighted trapezoidal filter is
\begin{equation}\label{eq:trap-filter}
  \frac1r\left(
  \frac{1+\cos(m\pi)}2+
  \sum_{j=1}^{r-1}\cos\frac{m\pi j}{r}
  \right)
  =
  \begin{cases}
  1,&2r\mid m,\\
  0,&2r\nmid m.
  \end{cases}
\end{equation}
Indeed, it is the real part of the periodic trapezoidal average
\[
  \frac1{2r}\sum_{\ell=0}^{2r-1}e^{i m\pi\ell/r}.
\]
Applying \eqref{eq:trap-filter} to the Fourier series of $\theta\mapsto L(\theta,\phi)$ gives
\begin{equation}\label{eq:L-average}
  \frac1r\left[
  \frac{L(0,\phi)+L(\pi,\phi)}2+
  \sum_{j=1}^{r-1}L(\pi j/r,\phi)
  \right]
  =
  \gamma(\phi)-\log2+
  \frac1r\log\left(1-e^{-2r\gamma(\phi)}\right).
\end{equation}

The preceding calculation is first made on $\phi\in[\varepsilon,\pi]$, where the Fourier series is uniformly convergent. Averaging \eqref{eq:gamma-integral} over the half-weighted $\theta$-grid and using \eqref{eq:L-average} gives the truncated identity on $[\varepsilon,\pi]$. It remains only to pass through the endpoint $\phi=0$. Since
\[
  \cos\phi=1-\frac{\phi^2}{2}+O(\phi^4),
  \qquad
  \gamma(\phi)=\arcosh\left(1+\frac{\phi^2}{2}+O(\phi^4)\right)
  =\phi+O(\phi^3),
\]
there is $\delta>0$ such that
\[
  \frac12\phi\le \gamma(\phi)\le 2\phi,
  \qquad 0<\phi\le\delta.
\]
Consequently, for fixed $r$,
\[
  \log\left(1-e^{-2r\gamma(\phi)}\right)
  =
  \log\phi+O_r(1),
  \qquad \phi\downarrow0.
\]
Equivalently, there is a constant $D_r>0$, depending only on $r$, such that
\[
  \left|\log\left(1-e^{-2r\gamma(\phi)}\right)\right|
  \le D_r(1+|\log\phi|),
  \qquad 0<\phi\le\delta,
\]
which is integrable at $0$.
On the left side, the only singular grid term is
\[
  L(0,\phi)=\log(1-\cos\phi)=2\log\phi+O(1);
\]
the remaining grid terms and $\gamma(\phi)-\log2$ are bounded near $0$. Hence both sides have at worst logarithmic endpoint singularities, with an integrable majorant independent of the cutoff $\varepsilon$. The contribution of $0<\phi<\varepsilon$ is $O_r(\varepsilon|\log\varepsilon|)$, and letting $\varepsilon\downarrow0$ gives \eqref{eq:trap-identity}.

It remains to prove monotonicity. For $z>0$ and $p>0$, set
\[
  h_p(z)=\frac1p\log(1-e^{-2pz}).
\]
Then
\[
  \frac{\partial}{\partial p}h_p(z)
  =
  -\frac1{p^2}\log(1-e^{-2pz})
  +
  \frac1p\frac{2z}{e^{2pz}-1}
  >0.
\]
Indeed, $0<1-e^{-2pz}<1$, so the logarithmic term is positive after the minus sign, and the second term is plainly positive. Since $\gamma(\phi)>0$ for $0<\phi\le\pi$, the already established identity \eqref{eq:trap-identity} gives, pointwise in $\phi$, $h_s(\gamma(\phi))\ge h_r(\gamma(\phi))$ whenever $s\ge r$. The endpoint bound above gives integrability, so no differentiation under the integral is needed. Therefore $T_s\ge T_r$, and
\[
  \frac{C_s}{s}+\frac{\alpha}{s}
  \ge
  \frac{C_r}{r}+\frac{\alpha}{r},
\]
which is equivalent to \eqref{eq:C-trapezoidal}.
\end{proof}

Multiplying \eqref{eq:trap-identity} by $r$ and using
$T_r=C_r/r+\alpha/r$ gives the equivalent error formula
\begin{equation}\label{eq:C-error}
  E_r:=rI-\alpha-C_r
  =
  -\frac1\pi
  \int_0^\pi
  \log\left(1-e^{-2r\gamma(\phi)}\right)\dd\phi.
\end{equation}
This is still part of the trapezoidal estimate, but it is written in $C_r$-notation for later use.
Since $\gamma$ is concave, $\gamma(0)=0$, and $\gamma(\pi)=c(1)$,
\[
  \gamma(\phi)\ge \frac{c(1)}{\pi}\phi.
\]
Consequently,
\begin{equation}\label{eq:C-error-bound}
  0\le E_r
  \le
  -\frac1\pi
  \int_0^\pi
  \log\left(1-e^{-2r c(1)\phi/\pi}\right)\dd\phi
  \le
  \frac{\pi^2}{12c(1)r}.
\end{equation}
Indeed, in the last integral the change of variables $u=2rc(1)\phi/\pi$ gives an upper bound
\[
  \frac{1}{2rc(1)}
  \int_0^\infty -\log(1-e^{-u})\dd u
  =
  \frac{\pi^2}{12c(1)r},
\]
using the standard identity
\[
  \int_0^\infty -\log(1-e^{-u})\dd u=\frac{\pi^2}{6}.
\]

\section{Proof of the Balancing Theorem and Optimality}\label{sec:proof-optimality}

\begin{proof}[Proof of \Cref{thm:balancing}, except optimality]
If $t=1$ in the notation of \Cref{thm:exact-gain}, the rectangles are identical. Assume $t>1$. Then $b>A$, and \Cref{thm:exact-gain} gives
\[
  \log\frac{\tau(a,b)}{\tau(A,B)}
  =
  (t-1)(AC_b-bC_A)+\Gamma
\]
with $\Gamma>0$. By \Cref{lem:trap}, applied with $r=A$ and $s=b$,
\[
  \frac{C_b}{b}-\frac{C_A}{A}
  \ge
  \frac{b-A}{Ab}\alpha.
\]
Hence
\[
  AC_b-bC_A
  =
  Ab\left(\frac{C_b}{b}-\frac{C_A}{A}\right)
  \ge
  \alpha(b-A).
\]
Therefore
\[
  \log\frac{\tau(a,b)}{\tau(A,B)}
  \ge
  \alpha(t-1)(b-A).
\]
Finally, since $a=tA$ and $B=tb$,
\[
  (t-1)(b-A)
  =
  (A+B)-(a+b).
\]
This proves the displayed inequality. If the rectangles are distinct, then $t>1$, and the strict inequality follows from $\Gamma>0$.
\end{proof}

For an $r\times s$ rectangle in the infinite grid, the number of outer boundary incidences is $\partial R_{r,s}=2r+2s$.
Thus \Cref{thm:balancing} can also be written as
\[
  \log\tau(a,b)-\log\tau(A,B)
  \ge
  \frac{\alpha}{2}
  \bigl(\partial R_{A,B}-\partial R_{a,b}\bigr).
\]
The side-length formulation avoids ambiguity about boundary conventions.

\begin{proposition}\label{prop:optimality}
The constant $\alpha=\arsinh(1)$ in \Cref{thm:balancing} is optimal.
\end{proposition}

\begin{proof}
Let
\[
  (A_n,B_n)=\bigl(n^2,(n+1)^2\bigr),
  \qquad
  (a_n,b_n)=\bigl(n(n+1),n(n+1)\bigr).
\]
Then $A_nB_n=a_nb_n$ and
\[
  (A_n+B_n)-(a_n+b_n)=1.
\]
Put $A=n^2$, $b=n(n+1)$, and $t_n=1+1/n$. The exact gain identity gives
\[
  L_n
  :=
  \log\frac{\tau(n(n+1),n(n+1))}
  {\tau(n^2,(n+1)^2)}
  =
  nC_b-(n+1)C_A+\Gamma_n.
\]
Use the $C_r$-error notation \eqref{eq:C-error}. Since $C_m=mI-\alpha-E_m$ and $nb=(n+1)A$,
\[
  nC_b-(n+1)C_A
  =
  \alpha+(n+1)E_A-nE_b.
\]
The monotonicity estimate \eqref{eq:C-trapezoidal}, equivalently $T_b\ge T_A$, gives $E_A/A\ge E_b/b$. Since $A=n^2$ and $b=n(n+1)$, this implies $(n+1)E_A\ge nE_b$. Therefore the term $(n+1)E_A-nE_b$ is nonnegative, and the error bound \eqref{eq:C-error-bound} gives
\[
  0\le (n+1)E_A-nE_b
  \le
  (n+1)E_A
  \le
  \frac{\pi^2(n+1)}{12c(1)n^2}
  \to0.
\]
Thus $nC_b-(n+1)C_A\to\alpha$.

It remains to show $\Gamma_n\to0$. Write $t_n=1+\delta_n$ with $\delta_n=1/n$. For $z>0$,
\[
  G_{1+\delta_n}(z)
  =
  \int_1^{1+\delta_n}\frac{2z}{e^{2vz}-1}\dd v
  \le
  \delta_n e^{-z},
\]
because
\[
  \frac{2z}{e^{2z}-1}
  =
  e^{-z}\frac{z}{\sinh z}
  \le e^{-z}.
\]
By concavity of $c$ and $c(0)=0$, we have $c(x)\ge c(1)x$. Hence, for the first sum in $\Gamma_n$,
\[
  A c(k/b)
  \ge
  n^2c(1)\frac{k}{n(n+1)}
  =
  c(1)\frac{n}{n+1}k
  \ge
  \frac{c(1)}2 k,
\]
and for the second sum,
\[
  b c(j/A)
  \ge
  n(n+1)c(1)\frac{j}{n^2}
  =
  c(1)\left(1+\frac1n\right)j
  \ge
  c(1)j.
\]
Therefore
\[
  |\Gamma_n|
  \le
  \frac1n\sum_{k\ge1}e^{-c(1)k/2}
  +
  \frac1n\sum_{j\ge1}e^{-c(1)j}
  =
  O(1/n),
\]
so $\Gamma_n\to0$. Hence $L_n\to\alpha$. Since the side-sum drop in this sequence is exactly $1$, no universal constant larger than $\alpha$ can satisfy the balancing inequality for all equal-area rectangular moves.
\end{proof}

For comparison, \Cref{app:concavity-constant} gives a direct concavity proof of the weaker balancing constant $\alpha_0=c(1)-\int_0^1 c(x)\,dx$; the proposition above shows why the sharp universal constant is instead $\alpha=c(1)/2=\arsinh(1)$.

\subsection*{A bounded-aspect asymptotic from the residual}

The balancing theorem uses only positivity of the residual. In bounded-aspect regimes, however, the same single-rectangle decomposition retains finer finite-size information. The following consequence recovers the classical area and boundary terms and the logarithmic corner correction, leaving only a bounded remainder.

\begin{proposition}\label{prop:bounded-aspect-residual}
Fix $M\ge1$. Uniformly for integers $2\le r\le s\le Mr$,
\[
  \log\tau(r,s)
  =
  Irs-\alpha(r+s)-\frac14\log(rs)+O_M(1)
  =
  \frac{4G_{\mathrm{Cat}}}{\pi}rs-\log(1+\sqrt2)(r+s)-\frac14\log(rs)+O_M(1).
\]
In particular, the classical square-grid area, boundary, and logarithmic corner terms are recovered:
\[
  \log\tau(n,n)
  =
  \frac{4G_{\mathrm{Cat}}}{\pi}n^2-2\log(1+\sqrt2)n-\frac12\log n+O(1),
\]
where $G_{\mathrm{Cat}}$ denotes Catalan's constant.
\end{proposition}

The coefficients in this proposition are not new as asymptotic constants. The area and boundary terms are classical in rectangular-grid spanning-tree asymptotics~\cite{Temperley1974,Wu1977}; the logarithmic corner correction appears in the finite-size correction literature for free-boundary rectangles. The point here is different: the same residual decomposition used to prove the exact finite balancing theorem also recovers these known terms, up to a bounded remainder, without invoking the full fixed-aspect expansion. The underlying finite product formula is the two-dimensional free-boundary case of the hypercubic-lattice formula of Tzeng and Wu~\cite{TzengWu2000}. Izmailian, Kenna, Guo, and Wu~\cite{IzmailianKennaGuoWu2014} compute the sharper fixed-aspect expansion, including the constant term; in their free-energy convention $F=-\log Z$, their bulk, surface, and corner contributions translate to the $Irs$, $-\alpha(r+s)$, and $-\frac14\log(rs)$ terms above. The square specialization is also consistent with the finite-size correction formulas of Izmailian and Kenna~\cite{IzmailianKenna2015}. We do not identify the remaining constant term here.

Even such a fixed-aspect refinement would not replace the finite balancing proof. If two rectangles have the same area, then the area and logarithmic corner terms cancel on subtracting the asymptotic formulas, and the boundary term favors the smaller side sum. The remaining fixed-aspect constant, however, need not be small, while some finite balancing moves have only a bounded side-sum drop. The balancing theorem therefore uses the exact gain identity and the finite positivity of the residual, rather than only this asymptotic information.

\begin{proof}
We use the same hyperbolic residual mechanism as in the gain identity, but for a single rectangle rather than for a ratio, so the numerator and denominator residual sums are kept separate.
The one-dimensional formula obtained from \eqref{eq:tau-q}, \eqref{eq:lambda-c}, and \Cref{lem:hyperbolic} gives, for $r\le s$,
\[
  \log\tau(r,s)
  =
  \sum_{k=1}^{r-1}
  \left[
  (s-1)c(k/r)
  +\log\left(1-e^{-2s c(k/r)}\right)
  -\log\left(1-e^{-2c(k/r)}\right)
  \right].
\]
Thus
\begin{equation}\label{eq:bounded-logtau-split}
  \log \tau(r,s)
  =
  (s-1)C_r
  +
  \sum_{k=1}^{r-1}\log\left(1-e^{-2s c(k/r)}\right)
  +
  \sum_{k=1}^{r-1}F(k/r),
\end{equation}
where
\[
  F(x)=-\log\left(1-e^{-2c(x)}\right).
\]
By \eqref{eq:C-error} and \eqref{eq:C-error-bound},
\[
  C_r=rI-\alpha-E_r,
  \qquad
  0\le E_r=O(1/r).
\]
Hence, uniformly for $s\le Mr$,
\begin{equation}\label{eq:bounded-main-term}
  (s-1)C_r
  =
  rsI-rI-\alpha s+O_M(1).
\end{equation}

Next, consider the numerator residual. By concavity of $c$ and $c(0)=0$, we have the explicit lower bound $c(x)\ge c(1)x$ on $[0,1]$. Taking $\beta=c(1)>0$, we have for $1\le k\le r-1$ that
\[
  sc(k/r)\ge rc(k/r)\ge\beta k.
\]
Every term in the numerator residual is nonpositive, and
\[
  \sum_{k=1}^{r-1}\log\left(1-e^{-2s c(k/r)}\right)
  \ge
  \sum_{k=1}^{\infty}\log\left(1-e^{-2\beta k}\right)
  >-\infty.
\]
Thus, uniformly for $s\ge r$,
\begin{equation}\label{eq:bounded-numerator-residual}
  \sum_{k=1}^{r-1}\log\left(1-e^{-2s c(k/r)}\right)=O(1).
\end{equation}

It remains to estimate the denominator residual. Since
\[
  c(x)=\pi x+O(x^3)
  \qquad (x\downarrow0),
\]
we have
\[
  F(x)=-\log x+h(x),
\]
where $h$ extends to a $C^1$ function on $[0,1]$, because $1-e^{-2c(x)}$ has a simple zero at $x=0$ with derivative $2\pi$. We use the following elementary logarithmic Riemann-sum estimate: if
\[
  f(x)=-\log x+h(x),
  \qquad h\in C^1[0,1],
\]
then
\begin{equation}\label{eq:log-riemann}
  \sum_{k=1}^{r-1}f(k/r)
  =
  r\int_0^1f(x)\dd x-\frac12\log r+O(1).
\end{equation}
Indeed,
\[
  \sum_{k=1}^{r-1}-\log(k/r)
  =
  (r-1)\log r-\log((r-1)!)
  =
  r-\frac12\log r+O(1)
\]
by Stirling's formula, while
\[
  \sum_{k=1}^{r-1}h(k/r)=r\int_0^1h(x)\dd x+O(1)
\]
by the ordinary Riemann-sum estimate for $C^1$ functions. Applying \eqref{eq:log-riemann} to $F$ gives
\begin{equation}\label{eq:bounded-denominator-residual}
  \sum_{k=1}^{r-1}F(k/r)=rJ-\frac12\log r+O(1),
  \qquad
  J:=\int_0^1F(x)\dd x.
\end{equation}

We now evaluate $J$. Put
\[
  \theta=\frac{\pi x}{2},
  \qquad
  u=\sin\theta,
  \qquad
  c(x)=2\arsinh u.
\]
If $a=\arsinh u$, then $c=2a$, and
\[
  1-e^{-2c}=1-e^{-4a}
  =
  4u\sqrt{1+u^2}\,e^{-c}.
\]
Therefore
\[
  F(x)=c(x)-\log\left(4\sin\theta\sqrt{1+\sin^2\theta}\right).
\]
Consequently,
\[
  J
  =
  I-\frac2\pi\int_0^{\pi/2}
  \log\left(4\sin\theta\sqrt{1+\sin^2\theta}\right)\dd\theta.
\]
Using
\[
  \int_0^{\pi/2}\log\sin\theta\dd\theta
  =
  -\frac{\pi}{2}\log2
\]
and
\[
  \int_0^{\pi/2}\log(1+\sin^2\theta)\dd\theta
  =
  \pi\log\frac{1+\sqrt2}{2},
\]
we obtain
\[
  \int_0^{\pi/2}
  \log\left(4\sin\theta\sqrt{1+\sin^2\theta}\right)\dd\theta
  =
  \frac{\pi}{2}\log(1+\sqrt2)
  =
  \frac{\pi}{2}\alpha.
\]
Thus
\begin{equation}\label{eq:F-integral}
  J=I-\alpha.
\end{equation}
Combining \eqref{eq:bounded-logtau-split}, \eqref{eq:bounded-main-term}, \eqref{eq:bounded-numerator-residual}, \eqref{eq:bounded-denominator-residual}, and \eqref{eq:F-integral} gives, uniformly for $2\le r\le s\le Mr$,
\[
  \log \tau(r,s)
  =
  rsI-rI-\alpha s+r(I-\alpha)-\frac12\log r+O_M(1)
  =
  Irs-\alpha(r+s)-\frac14\log(rs)+O_M(1),
\]
because $1\le s/r\le M$.
Finally,
\[
  I=\frac4\pi\int_0^{\pi/2}\arsinh(\sin\theta)\dd\theta.
\]
If $K(a)=\int_0^{\pi/2}\arsinh(a\sin\theta)\dd\theta$, then
\[
  K'(a)=\int_0^{\pi/2}\frac{\sin\theta}{\sqrt{1+a^2\sin^2\theta}}\dd\theta
  =
  \frac{\arctan a}{a}.
\]
Therefore $K(1)=\int_0^1(\arctan a)/a\dd a=G_{\mathrm{Cat}}$, giving the equivalent form.
\end{proof}

For a nontrivial balancing move $A\times B\to a\times b$, define its effective coefficient by
\[
  \alpha_{\mathrm{eff}}(A,B;a,b)
  =
  \frac{\log\tau(a,b)-\log\tau(A,B)}
  {(A+B)-(a+b)}.
\]
By \Cref{thm:balancing}, every such coefficient is at least $\alpha$, and \Cref{prop:optimality} shows that the infimum over all nontrivial equal-area rectangular balancing moves is $\alpha$. More generally, for any finite equal-area balancing sequence
\[
  R_0\to R_1\to\cdots\to R_m,
\]
the same lower bound holds after summing the one-step inequalities: the total log-gain is at least $\alpha$ times the total side-sum drop. For the direct move $1\times n^2\to n\times n$,
\[
  \alpha_{\square}(n)
  =
  \frac{\log\tau(n,n)}{(n-1)^2}
  =
  I+O(1/n),
\]
by \Cref{prop:bounded-aspect-residual}. Thus the same exact decomposition gives both the sharp worst-case balancing constant $\alpha$ and near-tight asymptotics for bounded-aspect rectangles through the residual; in particular, the shape-dependent coefficient for the path-to-square comparison tends to the square-lattice entropy $I$, even though the worst-case universal infimum remains $\alpha$.

\begin{corollary}\label{cor:square-max}
Among rectangular grid graphs with area $N$, the value $\tau(d,N/d)$ is strictly increasing as $d$ ranges over divisors $d\le\sqrt N$. Hence the unique maximizing rectangle, up to rotation, is
\[
  d_*\times \frac{N}{d_*},
  \qquad
  d_*=\max\{d:d\mid N,\ d\le\sqrt N\}.
\]
In particular, if $N=n^2$, the unique maximizing rectangle is $n\times n$.
\end{corollary}

\begin{proof}
Let $d_1<d_2\le\sqrt N$ be divisors of $N$. Set
\[
  A=d_1,\quad B=N/d_1,\quad a=d_2,\quad b=N/d_2.
\]
Then $AB=ab=N$ and $A<a\le b<B$. By \Cref{thm:balancing},
\[
  \log\tau(d_2,N/d_2)-\log\tau(d_1,N/d_1)>0.
\]
\end{proof}

\section{Discussion}\label{sec:discussion}

The argument uses the Cartesian-product spectrum of $P_r\square P_s$ in an essential way. After the one-dimensional factors cancel the Kirchhoff normalization, the remaining interaction product can be reorganized through the Chebyshev identity for $q_r$. The exact gain identity then separates the comparison into a main term controlled by a monotone trapezoidal identity for the special hyperbolic function $c$ and a residual term controlled by monotone matching. The optimal constant is $\alpha=\arsinh(1)=c(1)/2$.
The residual term is used only through positivity for the balancing theorem, but in bounded-aspect regimes it contains finer asymptotic information, as illustrated by the rectangular asymptotic above.

At the level of evaluation and decomposition, the hyperbolic reduction is broader than rectangles. If $H$ is a connected graph on $m$ vertices with positive Laplacian eigenvalues $\mu_1,\ldots,\mu_{m-1}$, then the same cancellation of the Kirchhoff normalization gives
\[
  \tau(P_r\square H)
  =
  \tau(H)\prod_{i=1}^{m-1}
  \frac{\sinh(rz_i)}{\sinh z_i},
  \qquad
  z_i=\arcosh(1+\mu_i/2).
\]
Thus the method gives one-dimensional evaluations not only for rectangles but also for cylindrical grids, obtained by taking $H$ to be a cycle graph. Periodic grids, or rectangular tori, similarly have separable spectra and analogous one-dimensional formulas. These extensions are computational rather than extremal: the balancing theorem uses additional structure special to path spectra, especially the monotonicity and concavity of $c(x)=\arcosh(2-\cos\pi x)$, together with the boundary scale measured by side-sum reduction. For tori there is no corresponding boundary term, and any analogue would require a different folded-spectrum comparison, likely involving aspect ratio rather than rectangular boundary reduction. Arbitrary induced grid subgraphs lack the product spectrum, so the method does not directly apply there; this does not rule out approximation schemes based on product-like decompositions or spectral comparison, but such estimates would require separate error control.

\Cref{conj:induced-grid} remains open. A simple extremal calculation points in the same direction but is too coarse to settle spanning-tree maximality: the square is already extremal for edge count and cycle rank.

Let $S\subset\Z^2$ have $|S|=n^2$, and suppose that the induced graph $G=\mathcal L[S]$ is connected. Let $E=|E(G)|$, and let
\[
  \rho(G)=|E(G)|-|V(G)|+1=E-n^2+1
\]
be its cycle rank. Then
\[
  E\le 2n(n-1),
  \qquad
  \rho(G)\le (n-1)^2,
\]
with equality only for the square $R_{n,n}$ up to translation.

Here is the short proof. Let $w$ and $h$ be the numbers of occupied columns and rows. Since $G$ is connected, its projections onto the horizontal and vertical axes are intervals, so the occupied bounding box contains $wh$ lattice points and $wh\ge n^2$. Let $\beta$ be the number of ordered pairs $(v,u)$ such that $v\in S$, $u\notin S$, and $u$ is a grid-neighbor of $v$. In each occupied row, the leftmost and rightmost occupied vertices have horizontal neighbors outside $S$, and similarly each occupied column contributes at least two vertical boundary incidences. Thus
\[
  \beta\ge 2w+2h\ge4\sqrt{wh}\ge4n.
\]
The identity $4n^2=2E+\beta$ follows by summing the four incident lattice directions over all vertices in $S$: internal edges are counted twice, while boundary incidences are counted once. Hence $E\le2n(n-1)$ and $\rho(G)\le(n-1)^2$. Equality forces $w=h=n$ and $wh=n^2$, so $S$ fills its $n\times n$ bounding box.

Thus the square is extremal for some coarse graph parameters. These parameters alone cannot prove spanning-tree maximality, since edge count and cycle rank do not determine the number of spanning trees. The rectangular family admits a stronger statement: exact finite-size comparisons, the optimal balancing constant, and bounded-aspect asymptotic information from the residual.

What is missing for arbitrary induced subgraphs is not merely a sharper perimeter estimate. The rectangular proof uses the product spectrum of $P_r\square P_s$ to turn the full spanning-tree ratio into a one-dimensional hyperbolic comparison, with a monotone main term and a positive matched residual. General induced subgraphs have no comparable separable spectrum, so the rectangular residual argument has no immediate analogue beyond product families. Thus \Cref{conj:induced-grid} remains a natural benchmark for understanding how geometry influences spanning-tree counts in grid graphs.

\bibliographystyle{alpha}
\bibliography{ref}

\appendix

\section{A Concavity Constant}\label{app:concavity-constant}

This appendix records the concavity-based constant used in an earlier version of the argument. The proof is included for intuition: a direct grid-insertion concavity estimate already implies a universal balancing inequality, although with a weaker balancing constant than the sharp one proved in the main text. The constant has the simple form
\[
  \alpha_0
  :=
  c(1)-\int_0^1 c(x)\dd x
  =
  c(1)-I
  =
  2\arsinh(1)-\frac{4G_{\mathrm{Cat}}}{\pi}
  \approx
  0.59650.
\]
The main theorem uses the sharper constant $\alpha=c(1)/2=\arsinh(1)$.

The elementary input is the following normalized-grid comparison for concave functions.

For a function $g:[0,1]\to\mathbb R$, define
\[
  M_r(g)=\frac1r\sum_{j=1}^{r-1}g(j/r).
\]

\begin{lemma}\label{lem:app-grid-insertion}
Let $g:[0,1]\to\mathbb R$ be concave and satisfy $g(0)=0$. For every $r\ge1$,
\[
  M_{r+1}(g)-M_r(g)
  \ge
  \frac{1}{r(r+1)}
  \left(g(1)-\int_0^1g(x)\dd x\right).
\]
Consequently, for $1\le r\le s$,
\[
  M_s(g)-M_r(g)
  \ge
  \frac{s-r}{rs}
  \left(g(1)-\int_0^1g(x)\dd x\right).
\]
\end{lemma}

\begin{proof}
Set $x_k=k/r$ and $y_j=j/(r+1)$. For $j=1,\ldots,r$, the point $y_j$ lies in $[x_{j-1},x_j]$, and
\[
  y_j
  =
  \frac{j}{r+1}x_{j-1}
  +
  \frac{r+1-j}{r+1}x_j.
\]
By concavity,
\[
  g(y_j)
  \ge
  \frac{j}{r+1}g(x_{j-1})
  +
  \frac{r+1-j}{r+1}g(x_j).
\]
Summing over $j=1,\ldots,r$ and using $g(0)=0$ gives
\[
  \sum_{j=1}^r g\left(\frac{j}{r+1}\right)
  \ge
  \frac{r+2}{r+1}\sum_{k=1}^{r-1}g\left(\frac{k}{r}\right)
  +
  \frac{1}{r+1}g(1).
\]
After dividing by $r+1$,
\[
  M_{r+1}(g)
  \ge
  \frac{r(r+2)}{(r+1)^2}M_r(g)
  +
  \frac{g(1)}{(r+1)^2}.
\]
Hence
\begin{equation}\label{eq:app-grid-insertion-step}
  M_{r+1}(g)-M_r(g)
  \ge
  \frac{g(1)-M_r(g)}{(r+1)^2}.
\end{equation}

Next we prove
\begin{equation}\label{eq:app-grid-insertion-integral}
  \int_0^1g(x)\dd x
  \ge
  \frac{1}{r+1}\sum_{j=1}^r g(j/r).
\end{equation}
Concavity implies that $g$ lies above its piecewise linear interpolation on the grid $0,1/r,\ldots,1$, so
\[
  \int_0^1g(x)\dd x
  \ge
  \frac1r\left(\sum_{j=1}^{r-1}g(j/r)+\frac12g(1)\right).
\]
Also, since $g(0)=0$, concavity gives
\[
  g(j/r)\ge \frac jr g(1),
  \qquad j=1,\ldots,r-1.
\]
Writing $S=\sum_{j=1}^{r-1}g(j/r)$, we get $S\ge (r-1)g(1)/2$. Therefore
\[
  \frac1r\left(S+\frac12g(1)\right)
  -
  \frac{S+g(1)}{r+1}
  =
  \frac{S-\frac{r-1}{2}g(1)}{r(r+1)}
  \ge 0,
\]
which proves \eqref{eq:app-grid-insertion-integral}. Rearranging \eqref{eq:app-grid-insertion-integral} gives
\[
  g(1)-M_r(g)
  \ge
  \frac{r+1}{r}
  \left(g(1)-\int_0^1g(x)\dd x\right).
\]
Combining this with \eqref{eq:app-grid-insertion-step} proves the one-step inequality.

Summing the one-step inequality from $r$ to $s-1$ gives
\[
  M_s(g)-M_r(g)
  \ge
  \left(\sum_{m=r}^{s-1}\frac{1}{m(m+1)}\right)
  \left(g(1)-\int_0^1g(x)\dd x\right),
\]
and the telescoping sum is $(s-r)/(rs)$.
\end{proof}

\begin{corollary}\label{cor:app-C-balancing}
For $1\le r\le s$,
\begin{equation}\label{eq:app-C-balancing}
  \frac{C_s}{s}-\frac{C_r}{r}
  \ge
  \frac{s-r}{rs}\alpha_0,
  \qquad
  \alpha_0=c(1)-I.
\end{equation}
\end{corollary}

\begin{proof}
Apply \Cref{lem:app-grid-insertion} to $g=c$. By \Cref{lem:c}, the function $c$ is concave and satisfies $c(0)=0$, and by definition $M_r(c)=C_r/r$.
\end{proof}

\begin{lemma}\label{lem:app-alpha-zero}
The comparison constant has the closed form
\[
  \alpha_0
  =
  2\arsinh(1)-\frac{4G_{\mathrm{Cat}}}{\pi}
  =
  \log(3+2\sqrt{2})-\frac{4G_{\mathrm{Cat}}}{\pi}
  \approx 0.59650,
\]
where
\[
  G_{\mathrm{Cat}}=\sum_{n\ge0}\frac{(-1)^n}{(2n+1)^2}
\]
is Catalan's constant.
\end{lemma}

\begin{proof}
First, $c(1)=2\arsinh(1)=\log(3+2\sqrt{2})$. Second,
\[
  I=\int_0^1 c(x)\dd x
  =
  \frac4\pi\int_0^{\pi/2}\arsinh(\sin\theta)\dd\theta.
\]
Let $K(u)=\int_0^{\pi/2}\arsinh(u\sin\theta)\dd\theta$. Then
\[
  K'(u)
  =
  \int_0^{\pi/2}
  \frac{\sin\theta}{\sqrt{1+u^2\sin^2\theta}}\dd\theta
  =
  \frac{\arctan u}{u},
\]
with limiting value $1$ at $u=0$. Therefore
\[
  K(1)=\int_0^1\frac{\arctan u}{u}\dd u
  =
  \sum_{n\ge0}\frac{(-1)^n}{(2n+1)^2}
  =
  G_{\mathrm{Cat}}.
\]
Thus $I=4G_{\mathrm{Cat}}/\pi$, and the formula for $\alpha_0=c(1)-I$ follows.
\end{proof}

\begin{proposition}\label{prop:app-old-balancing}
Under the hypotheses of \Cref{thm:balancing},
\[
  \log\tau(a,b)-\log\tau(A,B)
  \ge
  \alpha_0\bigl((A+B)-(a+b)\bigr).
\]
If the rectangles are distinct, the inequality is strict.
\end{proposition}

\begin{proof}
If $t=1$ in the notation of \Cref{thm:exact-gain}, the rectangles are identical. Assume $t>1$. Then \Cref{thm:exact-gain} gives
\[
  \log\frac{\tau(a,b)}{\tau(A,B)}
  =
  (t-1)(AC_b-bC_A)+\Gamma
\]
with $\Gamma>0$. Applying \eqref{eq:app-C-balancing} with $r=A$ and $s=b$ gives
\[
  AC_b-bC_A
  =
  Ab\left(\frac{C_b}{b}-\frac{C_A}{A}\right)
  \ge
  \alpha_0(b-A).
\]
Since $(t-1)(b-A)=(A+B)-(a+b)$, the claimed inequality follows. Strictness for distinct rectangles follows from $\Gamma>0$.
\end{proof}

\end{document}
